\chapter{Chapter04 — LMN}
%%% generated by the blueprint pipeline from FABL.Chapter04.LMN
\section{Overview}

Declarations of the Lean module \texttt{FABL.Chapter04.LMN}.

\section{Declarations}

\begin{definition}[Neutral gate value]
\lean{FABL.CircuitGate.neutral}
\label{FABL.CircuitGate.neutral}
\leanok
\uses{FABL.CircuitGate, FABL.Sign}
The neutral value of an AND gate is $-1$, and the neutral value of an OR gate is $1$.
\end{definition}

\begin{theorem}[Neutral value of an AND gate]
\lean{FABL.CircuitGate.neutral_and}
\label{FABL.CircuitGate.neutral_and}
\leanok
\uses{FABL.CircuitGate.neutral}
\difficulty{1}
For the sign-valued neutral map on circuit gates, the AND gate has neutral value $-1$.
\end{theorem}

\begin{theorem}[Neutral value of an OR gate]
\lean{FABL.CircuitGate.neutral_or}
\label{FABL.CircuitGate.neutral_or}
\leanok
\uses{FABL.CircuitGate.neutral}
\difficulty{1}
The neutral sign assigned to the logical OR gate is $1$.
\end{theorem}

\begin{definition}[Preimage of a coordinate set]
\lean{FABL.embeddingPreimage}
\label{FABL.embeddingPreimage}
\leanok
Let $m,n$ be natural numbers, let $e:\text{Fin}(m)\to\text{Fin}(n)$ be an embedding, and
let $I$ be a finite set of elements of $\text{Fin}(n)$. The preimage of $I$ under $e$ is
the set of all $i\in\text{Fin}(m)$ such that $e(i)\in I$.
\end{definition}

\begin{theorem}[Membership in an embedding preimage]
\lean{FABL.mem_embeddingPreimage}
\label{FABL.mem_embeddingPreimage}
\leanok
\uses{FABL.embeddingPreimage}
\difficulty{1}
Let $m,n$ be natural numbers, let $e:\bbr^{m}\to\bbr^{n}$ be an injective map between the
index sets $\bbr^{m}$ and $\bbr^{n}$, and let $A$ be a set of indices in $\bbr^{n}$. An
index $i$ belongs to the preimage of $A$ under $e$ if and only if $e(i)$ belongs to $A$.
\end{theorem}

\begin{theorem}[Gate evaluation under coordinate embedding]
\lean{FABL.CircuitGate.evalFinset_embeddingPreimage_extend}
\label{FABL.CircuitGate.evalFinset_embeddingPreimage_extend}
\leanok
\uses{FABL.CircuitGate, FABL.CircuitGate.evalFinset, FABL.CircuitGate.neutral, FABL.Sign, FABL.embeddingPreimage, FABL.mem_embeddingPreimage}
\difficulty{5}
Let $m,n$ be nonnegative integers, let $G$ be an AND or OR gate, let
$e:\\text{Fin}(m)\hookrightarrow\text{Fin}(n)$ be an injective coordinate map, let
$v:\text{Fin}(m)\to\{\pm1\}$ assign signs to the smaller coordinate set, and let
$S\subseteq\text{Fin}(n)$. Extend $v$ to the larger coordinate set by assigning the
neutral sign of $G$ outside the image of $e$. Then evaluating $G$ on the preimage of $S$
under $e$ with values $v$ gives the same sign as evaluating $G$ on $S$ with this
extended assignment.
\end{theorem}

\begin{definition}[Input cast for a circuit tail]
\lean{FABL.CircuitTail.castInput}
\label{FABL.CircuitTail.castInput}
\leanok
\uses{FABL.CircuitTail}
Given natural numbers $m$ and $k$, an equality $m=k$, and a circuit tail with $m$ input
coordinates, transporting it along that equality gives a circuit tail with $k$ input
coordinates.
\end{definition}

\begin{theorem}[Layer count under input reindexing]
\lean{FABL.CircuitTail.layerCount_castInput}
\label{FABL.CircuitTail.layerCount_castInput}
\leanok
\uses{FABL.CircuitTail, FABL.CircuitTail.castInput, FABL.CircuitTail.layerCount}
\difficulty{3}
Let $m,k$ be natural numbers with $m=k$, and let $T$ be a circuit tail with $m$ inputs.
Its number of layers is unchanged when its input indices are identified with those of
$k$.
\end{theorem}

\begin{theorem}[Internal node count under input reindexing]
\lean{FABL.CircuitTail.internalNodeCount_castInput}
\label{FABL.CircuitTail.internalNodeCount_castInput}
\leanok
\uses{FABL.CircuitTail, FABL.CircuitTail.castInput, FABL.CircuitTail.internalNodeCount}
\difficulty{3}
Let $m,k$ be natural numbers with $m=k$, and let $T$ be a circuit tail with $m$ inputs.
The number of internal nodes of $T$ is unchanged when its input indices are identified
with those of $k$.
\end{theorem}

\begin{theorem}[Evaluation under input reindexing]
\lean{FABL.CircuitTail.eval_castInput}
\label{FABL.CircuitTail.eval_castInput}
\leanok
\uses{FABL.CircuitGate, FABL.CircuitTail, FABL.CircuitTail.castInput, FABL.CircuitTail.eval, FABL.Sign}
\difficulty{3}
Let $m,k$ be natural numbers with $m=k$, let $T$ be a circuit tail with $m$ inputs, let
$g$ be an AND or OR gate, and let $v:\text{Fin}(k)\to\text{Sign}$ assign a sign to each of
the $k$ inputs. Evaluating $T$ after identifying its input indices with those of $k$
gives the same sign as evaluating $T$ on the values reindexed along that identification.
\end{theorem}

\begin{definition}[Pulling inputs through an embedding]
\lean{FABL.CircuitTail.pullInputs}
\label{FABL.CircuitTail.pullInputs}
\leanok
\uses{FABL.CircuitTail, FABL.CircuitTail.castInput, FABL.embeddingPreimage}
Given a circuit tail with $m$ input coordinates and an injective map from $k$
coordinates into those $m$ coordinates, its pullback is a circuit tail on the $k$
coordinates: each layer's gate input sets, including the output layer's, are replaced by
their preimages under the map.
\end{definition}

\begin{theorem}[Layer count under input pullback]
\lean{FABL.CircuitTail.layerCount_pullInputs}
\label{FABL.CircuitTail.layerCount_pullInputs}
\leanok
\uses{FABL.CircuitTail, FABL.CircuitTail.layerCount, FABL.CircuitTail.pullInputs}
\difficulty{4}
Let $m,k$ be natural numbers, let $T$ be a circuit tail with $m$ input indices, and let
$e:\text{Fin}(k)\to\text{Fin}(m)$ be an injective map. Pulling each layer's input sets back
along $e$ leaves the layer count unchanged.
\end{theorem}

\begin{theorem}[Internal node count after pulling back inputs]
\lean{FABL.CircuitTail.internalNodeCount_pullInputs}
\label{FABL.CircuitTail.internalNodeCount_pullInputs}
\leanok
\uses{FABL.CircuitTail, FABL.CircuitTail.internalNodeCount, FABL.CircuitTail.pullInputs}
\difficulty{4}
Let $T$ be a circuit tail with $m$ inputs, and let $e:\text{Fin}(k)\to\text{Fin}(m)$ be
injective. The number of internal nodes is unchanged when the circuit inputs are pulled
back along $e$.
\end{theorem}

\begin{theorem}[Evaluation after pulling back inputs]
\lean{FABL.CircuitTail.eval_pullInputs}
\label{FABL.CircuitTail.eval_pullInputs}
\leanok
\uses{FABL.CircuitGate, FABL.CircuitGate.evalFinset_embeddingPreimage_extend, FABL.CircuitGate.neutral, FABL.CircuitTail, FABL.CircuitTail.eval, FABL.CircuitTail.eval_castInput, FABL.CircuitTail.pullInputs, FABL.Sign}
\difficulty{4}
Let $T$ be a circuit tail with $m$ inputs, let $g$ be an AND or OR gate, let
$e:\text{Fin}(k)\to\text{Fin}(m)$ be an injective map, and assign a sign to each of the $k$
inputs. Pulling the circuit inputs back along $e$ and evaluating gives the same sign as
evaluating the original tail after assigning the given values on the image of $e$ and
assigning $-1$ outside that image for an AND gate, or $1$ outside it for an OR gate.
\end{theorem}

\begin{theorem}[Uniform probability of a false clause]
\lean{FABL.DepthCircuit.uniformProbability_clauseEval_one}
\label{FABL.DepthCircuit.uniformProbability_clauseEval_one}
\leanok
\uses{FABL.CNFFormula.clauseEval, FABL.CNFFormula.clauseEval_neg_iff_termEval, FABL.DNFTerm, FABL.DNFTerm.eval, FABL.DNFTerm.width, FABL.uniformProbability, FABL.uniformProbability_DNFTerm_eval_neg_one}
\difficulty{5}
For every term on the $n$-dimensional sign cube, the uniform probability that its clause
evaluation is $1$ equals $2^{-w}$, where $w$ is the number of literals in the term.
\end{theorem}

\begin{theorem}[First-layer gate output probability]
\lean{FABL.DepthCircuit.uniformProbability_evalLayer1_ne_neutral}
\label{FABL.DepthCircuit.uniformProbability_evalLayer1_ne_neutral}
\leanok
\uses{FABL.CNFFormula.clauseEval, FABL.CircuitGate.dual, FABL.CircuitGate.dual_and, FABL.CircuitGate.dual_or, FABL.CircuitGate.evalTerm, FABL.CircuitGate.neutral, FABL.CircuitGate.neutral_and, FABL.CircuitGate.neutral_or, FABL.DNFTerm.eval, FABL.DNFTerm.width, FABL.DepthCircuit, FABL.DepthCircuit.evalLayer1, FABL.DepthCircuit.uniformProbability_clauseEval_one, FABL.Sign, FABL.uniformProbability, FABL.uniformProbability_DNFTerm_eval_neg_one}
\difficulty{4}
For every depth circuit on the $n$-dimensional sign cube and every gate in its first
layer, the uniform probability that the gate's output differs from the neutral value of
the dual gate equals $2^{-w}$, where $w$ is the width of that gate's term.
\end{theorem}

\begin{definition}[Retained first-layer indices]
\lean{FABL.DepthCircuit.retainedLayer1Indices}
\statementsource{Analysis-of-Boolean-Functions-by-Ryan-ODonnell}{pdf-ffbe9c7b3306-p106-b005}
\label{FABL.DepthCircuit.retainedLayer1Indices}
\leanok
\uses{FABL.DNFTerm.width, FABL.DepthCircuit}
Let $C$ be a depth circuit on $n$ variables and let $w$ be a natural number. The
retained first-layer indices are the indices of those first-layer gates whose terms have
width at most $w$.
\end{definition}

\begin{definition}[Embedding of retained layer-one gates]
\lean{FABL.DepthCircuit.retainedLayer1Embedding}
\label{FABL.DepthCircuit.retainedLayer1Embedding}
\leanok
\uses{FABL.DepthCircuit, FABL.DepthCircuit.retainedLayer1Indices}
Let $C$ be a depth circuit and let $w$ be a natural number. The layer-one gates of $C$
whose terms have width at most $w$ form a finite set of indices. The definition gives an
injective map from an ordinal whose size is the number of these indices into the ordinal
indexing all layer-one gates, with image exactly the set of retained indices.
\end{definition}

\begin{theorem}[Retained-index embedding inverse law]
\lean{FABL.DepthCircuit.retainedLayer1Embedding_equivFin}
\label{FABL.DepthCircuit.retainedLayer1Embedding_equivFin}
\leanok
\uses{FABL.DepthCircuit, FABL.DepthCircuit.retainedLayer1Embedding, FABL.DepthCircuit.retainedLayer1Indices}
\difficulty{1}
Let $C$ be a depth circuit, let $w$ be a natural number, and let $i$ be a retained
first-layer index, meaning that the corresponding term has width at most $w$. Applying
the retained-index embedding to the finite-set equivalence's coordinate for $i$ returns
$i$.
\end{theorem}

\begin{theorem}[Image lies among retained indices]
\lean{FABL.DepthCircuit.retainedLayer1Embedding_mem}
\label{FABL.DepthCircuit.retainedLayer1Embedding_mem}
\leanok
\uses{FABL.DepthCircuit, FABL.DepthCircuit.retainedLayer1Embedding, FABL.DepthCircuit.retainedLayer1Indices}
\difficulty{3}
Let $C$ be a depth circuit, let $w$ be a natural number, and let $i$ index one of the
positions in the list of retained first-layer indices. Its image under the
retained-index embedding is itself a retained first-layer index.
\end{theorem}

\begin{definition}[Retained first-layer terms]
\lean{FABL.DepthCircuit.retainedLayer1}
\statementsource{Analysis-of-Boolean-Functions-by-Ryan-ODonnell}{pdf-ffbe9c7b3306-p106-b005}
\label{FABL.DepthCircuit.retainedLayer1}
\leanok
\uses{FABL.DNFTerm, FABL.DepthCircuit, FABL.DepthCircuit.retainedLayer1Embedding, FABL.DepthCircuit.retainedLayer1Indices}
Let $C$ be a depth circuit on $n$ variables and let $w$ be a natural number. The
retained first-layer terms are the list of terms at the retained first-layer indices, in
the order induced by those indices.
\end{definition}

\begin{theorem}[Length of the retained first layer]
\lean{FABL.DepthCircuit.length_retainedLayer1}
\label{FABL.DepthCircuit.length_retainedLayer1}
\leanok
\uses{FABL.DepthCircuit, FABL.DepthCircuit.retainedLayer1, FABL.DepthCircuit.retainedLayer1Indices}
\difficulty{1}
Let $C$ be a depth circuit and let $w$ be a natural number. The number of terms in the
list obtained by retaining exactly those first-layer terms of width at most $w$ equals
the number of retained first-layer indices.
\end{theorem}

\begin{theorem}[Indexing the retained first-layer terms]
\lean{FABL.DepthCircuit.get_retainedLayer1}
\label{FABL.DepthCircuit.get_retainedLayer1}
\leanok
\uses{FABL.DepthCircuit, FABL.DepthCircuit.length_retainedLayer1, FABL.DepthCircuit.retainedLayer1, FABL.DepthCircuit.retainedLayer1Embedding}
\difficulty{3}
Let $C$ be a layered circuit, let $w$ be a width threshold, and let $i$ index the list
of first-layer terms retained by the width truncation. The term at position $i$ in that
list equals the original first-layer term at the retained index corresponding to $i$.
\end{theorem}

\begin{definition}[Width truncation of a circuit]
\lean{FABL.DepthCircuit.truncateWidth}
\statementsource{Analysis-of-Boolean-Functions-by-Ryan-ODonnell}{pdf-ffbe9c7b3306-p106-b005}
\label{FABL.DepthCircuit.truncateWidth}
\leanok
\uses{FABL.CircuitTail.castInput, FABL.CircuitTail.pullInputs, FABL.DepthCircuit, FABL.DepthCircuit.length_retainedLayer1, FABL.DepthCircuit.retainedLayer1, FABL.DepthCircuit.retainedLayer1Embedding}
Let $C$ be a depth circuit on $n$ variables and let $w$ be a natural number. Its width
truncation retains from the first layer just the terms of width at most $w$, preserves
the first-layer gate type, and adjusts the remaining circuit layers to receive inputs
from the retained terms.
\end{definition}

\begin{theorem}[First-layer gate under width truncation]
\lean{FABL.DepthCircuit.layer1Gate_truncateWidth}
\label{FABL.DepthCircuit.layer1Gate_truncateWidth}
\leanok
\uses{FABL.DepthCircuit, FABL.DepthCircuit.truncateWidth}
\difficulty{1}
For any depth circuit $C$ and nonnegative integer $w$, the first-layer gate type of the
circuit obtained by truncating $C$ to width $w$ is the same as that of $C$.
\end{theorem}

\begin{theorem}[Width truncation preserves depth]
\lean{FABL.DepthCircuit.depth_truncateWidth}
\label{FABL.DepthCircuit.depth_truncateWidth}
\leanok
\uses{FABL.CircuitTail.layerCount_pullInputs, FABL.DepthCircuit, FABL.DepthCircuit.depth, FABL.DepthCircuit.retainedLayer1Embedding, FABL.DepthCircuit.truncateWidth}
\difficulty{1}
For every natural number $n$, every circuit on $n$ variables, and every natural number
width bound $w$, deleting first-layer gates whose width exceeds $w$ and pulling the
later circuit inputs back to the retained gates leaves the circuit depth unchanged.
\end{theorem}

\begin{theorem}[Size after width truncation]
\lean{FABL.DepthCircuit.size_truncateWidth_le}
\label{FABL.DepthCircuit.size_truncateWidth_le}
\leanok
\uses{FABL.CircuitTail.internalNodeCount, FABL.CircuitTail.internalNodeCount_castInput, FABL.CircuitTail.internalNodeCount_pullInputs, FABL.DepthCircuit, FABL.DepthCircuit.retainedLayer1Indices, FABL.DepthCircuit.size, FABL.DepthCircuit.truncateWidth}
\difficulty{3}
For every depth circuit and every natural-number width bound, the size of the circuit
obtained by retaining only first-layer terms of width at most that bound and restricting
the later layers accordingly is at most the size of the original circuit.
\end{theorem}

\begin{theorem}[Width after truncation]
\lean{FABL.DepthCircuit.width_truncateWidth_le}
\proofsource{Analysis-of-Boolean-Functions-by-Ryan-ODonnell}{pdf-ffbe9c7b3306-p106-b005}
\label{FABL.DepthCircuit.width_truncateWidth_le}
\leanok
\uses{FABL.DNFFormula.width_le_iff, FABL.DepthCircuit, FABL.DepthCircuit.get_retainedLayer1, FABL.DepthCircuit.length_retainedLayer1, FABL.DepthCircuit.retainedLayer1Embedding_mem, FABL.DepthCircuit.truncateWidth, FABL.DepthCircuit.width}
\difficulty{4}
Let $n$ and $w$ be nonnegative integers, and let $C$ be a depth circuit on $n$
coordinates. Form a circuit by retaining exactly those first-layer terms of $C$ whose
widths are at most $w$, and pulling the later circuit layers back to the retained gates.
The resulting circuit has width at most $w$.
\end{theorem}

\begin{theorem}[Evaluation after width truncation]
\lean{FABL.DepthCircuit.eval_truncateWidth}
\label{FABL.DepthCircuit.eval_truncateWidth}
\leanok
\uses{FABL.CircuitGate.dual, FABL.CircuitGate.evalTerm, FABL.CircuitGate.neutral, FABL.CircuitTail.castInput, FABL.CircuitTail.eval, FABL.CircuitTail.eval_castInput, FABL.CircuitTail.eval_pullInputs, FABL.CircuitTail.pullInputs, FABL.DepthCircuit, FABL.DepthCircuit.eval, FABL.DepthCircuit.evalLayer1, FABL.DepthCircuit.get_retainedLayer1, FABL.DepthCircuit.length_retainedLayer1, FABL.DepthCircuit.retainedLayer1, FABL.DepthCircuit.retainedLayer1Embedding, FABL.DepthCircuit.truncateWidth}
\difficulty{3}
Let $n$ be a natural number, $C$ a circuit on $n$ variables, $w$ a natural number, and
$x\in\{-1,1\}^n$. Retain precisely the first-layer gates of width at most $w$. The
truncated circuit's value at $x$ equals the value obtained by evaluating the original
later layers with each retained gate's value at $x$ on its corresponding input and each
omitted gate assigned the neutral sign for the dual of the first-layer gate type.
\end{theorem}

\begin{theorem}[Evaluation after width truncation]
\lean{FABL.DepthCircuit.eval_truncateWidth_eq_of_deleted_neutral}
\label{FABL.DepthCircuit.eval_truncateWidth_eq_of_deleted_neutral}
\leanok
\uses{FABL.CircuitGate.dual, FABL.CircuitGate.neutral, FABL.DepthCircuit, FABL.DepthCircuit.eval, FABL.DepthCircuit.evalLayer1, FABL.DepthCircuit.eval_truncateWidth, FABL.DepthCircuit.retainedLayer1Embedding, FABL.DepthCircuit.retainedLayer1Embedding_equivFin, FABL.DepthCircuit.retainedLayer1Indices, FABL.DepthCircuit.truncateWidth}
\difficulty{4}
Let $C$ be a finite-depth circuit with first-layer gate type $g$, let $w$ be a width
threshold, and let $x\in\{−1,1\}^n$. If every first-layer gate omitted by truncating to
gates of width at most $w$ evaluates at $x$ to the neutral value of the dual gate $g^*$
(that is, $-1$ for OR and $1$ for AND), then the truncated circuit and $C$ have the same
value at $x$.
\end{theorem}

\begin{theorem}[A deleted gate witnesses changed evaluation]
\lean{FABL.DepthCircuit.exists_deletedLayer1_of_eval_ne_truncateWidth}
\label{FABL.DepthCircuit.exists_deletedLayer1_of_eval_ne_truncateWidth}
\leanok
\uses{FABL.CircuitGate.dual, FABL.CircuitGate.neutral, FABL.DepthCircuit, FABL.DepthCircuit.eval, FABL.DepthCircuit.evalLayer1, FABL.DepthCircuit.eval_truncateWidth_eq_of_deleted_neutral, FABL.DepthCircuit.retainedLayer1Indices, FABL.DepthCircuit.truncateWidth}
\difficulty{3}
Let $C$ be a layered circuit, let $w$ be a width threshold, and let $x\in\{−1,1\}^n$. If
$C$ and its truncation retaining only first-layer gates of width at most $w$ have
different values at $x$, then some omitted first-layer gate evaluates at $x$ to a value
other than the neutral value of the dual of the first-layer gate type.
\end{theorem}

\begin{theorem}[Relative hamming dist truncate width le]
\lean{FABL.DepthCircuit.relativeHammingDist_truncateWidth_le}
\statementsource{Analysis-of-Boolean-Functions-by-Ryan-ODonnell}{pdf-ffbe9c7b3306-p106-b005}
\label{FABL.DepthCircuit.relativeHammingDist_truncateWidth_le}
\leanok
\uses{FABL.CircuitGate.dual, FABL.CircuitGate.neutral, FABL.CircuitTail.internalNodeCount, FABL.DNFTerm.width, FABL.DepthCircuit, FABL.DepthCircuit.eval, FABL.DepthCircuit.evalLayer1, FABL.DepthCircuit.exists_deletedLayer1_of_eval_ne_truncateWidth, FABL.DepthCircuit.retainedLayer1Indices, FABL.DepthCircuit.size, FABL.DepthCircuit.toBooleanFunction, FABL.DepthCircuit.truncateWidth, FABL.DepthCircuit.uniformProbability_evalLayer1_ne_neutral, FABL.relativeHammingDist, FABL.uniformProbability, FABL.uniformProbability_ne_eq_relativeHammingDist}
\difficulty{5}
% dataset-pipeline: stub -- prose pending build_informal_statements.py
Statement pending regeneration.
\end{theorem}

\begin{theorem}[Relative hamming dist truncate width cutoff le]
\lean{FABL.DepthCircuit.relativeHammingDist_truncateWidth_cutoff_le}
\statementsource{Analysis-of-Boolean-Functions-by-Ryan-ODonnell}{pdf-ffbe9c7b3306-p106-b005}
\label{FABL.DepthCircuit.relativeHammingDist_truncateWidth_cutoff_le}
\leanok
\uses{FABL.DepthCircuit, FABL.DepthCircuit.relativeHammingDist_truncateWidth_le, FABL.DepthCircuit.size, FABL.DepthCircuit.toBooleanFunction, FABL.DepthCircuit.truncateWidth, FABL.dnfWidthTruncationCutoff, FABL.mul_inv_two_pow_dnfWidthTruncationCutoff_le, FABL.relativeHammingDist}
\difficulty{4}
% dataset-pipeline: stub -- prose pending build_informal_statements.py
Statement pending regeneration.
\end{theorem}

\begin{theorem}[Width-truncated circuit representation]
\lean{FABL.DepthCircuit.hasDepthCircuit_truncateWidth}
\label{FABL.DepthCircuit.hasDepthCircuit_truncateWidth}
\leanok
\uses{FABL.DepthCircuit, FABL.DepthCircuit.HasDepthCircuit, FABL.DepthCircuit.depth, FABL.DepthCircuit.depth_truncateWidth, FABL.DepthCircuit.size, FABL.DepthCircuit.size_truncateWidth_le, FABL.DepthCircuit.toBooleanFunction, FABL.DepthCircuit.truncateWidth, FABL.DepthCircuit.width_truncateWidth_le}
\difficulty{1}
Let $n$ be a nonnegative integer, let $C$ be a depth circuit on the $n$-dimensional sign
cube, and let $w$ be a nonnegative integer. The function computed by the circuit
obtained from $C$ by truncating its first-layer terms to width at most $w$ has a
depth-$\mathrm{depth}(C)$ representation of size at most $\mathrm{size}(C)$ and width at
most $w$.
\end{theorem}

\begin{definition}[Switching layer rate]
\lean{FABL.switchingLayerRate}
\label{FABL.switchingLayerRate}
\leanok
For every natural number $w$, define the switching layer rate to be the real number
$1/(10w)$, with $w$ viewed as a real number in the denominator.
\end{definition}

\begin{definition}[Circuit compression rate]
\lean{FABL.circuitCompressionRate}
\label{FABL.circuitCompressionRate}
\leanok
\uses{FABL.switchingLayerRate}
For natural numbers $w$, $\ell$, and $r$, the circuit compression rate is the switching
layer rate at $w$ multiplied by the $r$th power of the switching layer rate at $\ell$.
\end{definition}

\begin{theorem}[Nonnegativity of the switching layer rate]
\lean{FABL.switchingLayerRate_nonneg}
\label{FABL.switchingLayerRate_nonneg}
\leanok
\uses{FABL.switchingLayerRate}
\difficulty{1}
For every natural number $w$, the switching layer rate $1/(10w)$ is nonnegative.
\end{theorem}

\begin{theorem}[Upper bound for the switching layer rate]
\lean{FABL.switchingLayerRate_le_one}
\label{FABL.switchingLayerRate_le_one}
\leanok
\uses{FABL.switchingLayerRate}
\difficulty{4}
For every positive integer $w$, the switching layer rate $1/(10w)$ is at most $1$.
\end{theorem}

\begin{theorem}[Nonnegativity of the circuit compression rate]
\lean{FABL.circuitCompressionRate_nonneg}
\label{FABL.circuitCompressionRate_nonneg}
\leanok
\uses{FABL.circuitCompressionRate, FABL.switchingLayerRate_nonneg}
\difficulty{1}
For all natural numbers $w,\ell,r$, define $a_j=1/(10j)$ for positive $j$ and set
$a_0=0$, and let $R(w,\ell,r)=a_w a_\ell^{r}$. Then $R(w,\ell,r)\geq 0$.
\end{theorem}

\begin{theorem}[Upper bound for the circuit compression rate]
\lean{FABL.circuitCompressionRate_le_one}
\label{FABL.circuitCompressionRate_le_one}
\leanok
\uses{FABL.circuitCompressionRate, FABL.switchingLayerRate_le_one, FABL.switchingLayerRate_nonneg}
\difficulty{3}
For natural numbers $w,\ell,r$ with $w>0$ and $\ell>0$, the circuit compression rate
$\frac{1}{10w}\left(\frac{1}{10\ell}\right)^r$ is at most $1$.
\end{theorem}

\begin{theorem}[Splitting the circuit compression rate]
\lean{FABL.circuitCompressionRate_split}
\label{FABL.circuitCompressionRate_split}
\leanok
\uses{FABL.circuitCompressionRate, FABL.switchingLayerRate}
\difficulty{4}
For natural numbers $w$, $\ell$, and $r$ with $1 \le r$, define $s(k)=1/(10k)$ for each
natural number $k$, viewing $k$ as a real number, and define $C(a,b,q)=s(a)s(b)^q$. Then
$s(w)C(\ell,\ell,r-1)=C(w,\ell,r)$.
\end{theorem}

\begin{definition}[Half switching bound]
\lean{FABL.HasHalfSwitchingBound}
\statementsource{lmn-constant-depth-fourier}{pdf-63759bd909d5-p005-b004}
\statementsource{Analysis-of-Boolean-Functions-by-Ryan-ODonnell}{pdf-ffbe9c7b3306-p100-b005}
\label{FABL.HasHalfSwitchingBound}
\leanok
\uses{FABL.BooleanFunction, FABL.HasCNFWidthLE, FABL.HasDNFWidthLE, FABL.switchingFailureProbability, FABL.switchingLayerRate}
For every Boolean function $f$ on $N$ bits and every positive integer $w$, suppose $f$
has a disjunctive normal form or a conjunctive normal form of width at most $w$. Then,
for every nonnegative integer $k$, the probability that $f$ fails under switching at
rate determined by $w$ is at most $2^{-k}$.
\end{definition}

\begin{theorem}[Restriction commutes with finite gate evaluation]
\lean{FABL.CircuitGate.reindexedSignRestriction_evalFinsetFunction}
\label{FABL.CircuitGate.reindexedSignRestriction_evalFinsetFunction}
\leanok
\uses{FABL.BooleanFunction, FABL.CircuitGate, FABL.CircuitGate.evalFinsetFunction, FABL.FixedSignCube, FABL.reindexedSignRestriction}
\difficulty{1}
Let $n,m$ be natural numbers, choose an AND or OR gate, a finite set of input indices
among $m$ indices, and a Boolean function on the $n$-dimensional sign cube for each
input index. For any set $J$ of coordinates and any fixed signs outside $J$, restricting
the function obtained by applying the gate to the selected input functions gives the
same function as first restricting each input function and then applying the gate to the
same selected inputs.
\end{theorem}

\begin{theorem}[Width from restricted decision tree depth]
\lean{FABL.CircuitGate.hasWidthLE_reindexedSignRestriction_of_depth_le}
\label{FABL.CircuitGate.hasWidthLE_reindexedSignRestriction_of_depth_le}
\leanok
\uses{FABL.BooleanFunction, FABL.CircuitGate, FABL.CircuitGate.HasWidthLE, FABL.CircuitGate.hasWidthLE_decisionTreeDepth, FABL.CircuitGate.hasWidthLE_mono, FABL.FixedSignCube, FABL.F₂DecisionTree.decisionTreeDepth, FABL.binaryOfBooleanFunction, FABL.reindexedSignRestriction, FABL.restrictedDecisionTreeDepth, FABL.restrictedDecisionTreeDepth_eq_reindexedSignRestriction}
\difficulty{4}
Let $n$ be a natural number, let $f$ be a function from the $n$-dimensional sign cube to
$\\{−1,1\\}$, let $J$ be a set of coordinates, and fix signs on the coordinates outside
$J$. If the resulting restriction of $f$, reindexed on the $|J|$ free coordinates, has
decision-tree depth at most $\\ell$, then for either gate choice it has a representation
of width at most $\\ell$: a CNF when the gate is AND, or a DNF when the gate is OR.
\end{theorem}

\begin{theorem}[Restriction commutes with circuit evaluation]
\lean{FABL.CircuitTail.reindexedSignRestriction_evalFunction}
\label{FABL.CircuitTail.reindexedSignRestriction_evalFunction}
\leanok
\uses{FABL.BooleanFunction, FABL.CircuitGate, FABL.CircuitTail, FABL.CircuitTail.evalFunction, FABL.FixedSignCube, FABL.reindexedSignRestriction}
\difficulty{1}
Let $n,m\in\mathbb{N}$, let $T$ be a circuit tail with $m$ inputs, let $g$ be an AND or
OR gate, and assign a sign-valued Boolean function on the $n$-dimensional sign cube to
each input of $T$. For a set $J$ of coordinates and a fixed sign assignment $z$ on its
complement, reindex the free coordinates of each restricted function by the $|J|$
positions of $J$. Then restricting the function computed by $T$ gives exactly the
function computed by $T$ from the individually restricted input functions, using the
same gate $g$.
\end{theorem}

\begin{theorem}[Switching bound for circuit evaluation]
\lean{FABL.CircuitTail.switchingFailureProbability_evalFunction_le_of_halfSwitching}
\proofsource{Analysis-of-Boolean-Functions-by-Ryan-ODonnell}{
  pdf-ffbe9c7b3306-p104-b004,
  pdf-ffbe9c7b3306-p104-b005,
  pdf-ffbe9c7b3306-p105-b001,
  pdf-ffbe9c7b3306-p106-b002
}
\label{FABL.CircuitTail.switchingFailureProbability_evalFunction_le_of_halfSwitching}
\leanok
\uses{FABL.BooleanFunction, FABL.CircuitGate, FABL.CircuitGate.HasWidthLE, FABL.CircuitGate.dual, FABL.CircuitGate.evalFinsetFunction, FABL.CircuitGate.hasDNFWidthLE_or_hasCNFWidthLE, FABL.CircuitGate.hasWidthLE_evalFinsetFunction, FABL.CircuitGate.hasWidthLE_reindexedSignRestriction_of_depth_le, FABL.CircuitTail, FABL.CircuitTail.evalFunction, FABL.CircuitTail.evalFunction_output, FABL.CircuitTail.internalNodeCount, FABL.CircuitTail.layerCount, FABL.CircuitTail.one_le_layerCount, FABL.FixedSignCube, FABL.HasHalfSwitchingBound, FABL.circuitCompressionRate, FABL.circuitCompressionRate_le_one, FABL.circuitCompressionRate_nonneg, FABL.circuitCompressionRate_split, FABL.expectRandomRestriction, FABL.expectRandomRestriction_add, FABL.expectRandomRestriction_const, FABL.expectRandomRestriction_mono, FABL.expectRandomRestriction_sum, FABL.reindexedSignRestriction, FABL.restrictedDecisionTreeDepth, FABL.switchingFailureIndicator, FABL.switchingFailureIndicator_nonneg, FABL.switchingFailureProbability, FABL.switchingFailureProbability_le_one, FABL.switchingFailureProbability_mul, FABL.switchingLayerRate, FABL.switchingLayerRate_le_one, FABL.switchingLayerRate_nonneg}
\difficulty{6}
Assume that every Boolean function having a disjunctive or conjunctive normal form of
width at most a positive integer $w$ has switching failure probability at rate $1/(10w)$
bounded by $2^{-k}$ for every natural number $k$. Let $T$ be a circuit tail with $m$
inputs, let $g$ be an AND or OR gate, and let each of its input functions have a normal
form appropriate to $g$ of width at most $w$. For positive integers $w$ and $\ell$, the
switching failure probability of the function computed by $T$ at rate
$\frac{1}{10w}\left(\frac{1}{10\ell}\right)^{L-1}$ is at most $N2^{-\ell}+2^{-k}$ for
every $k$, where $L$ is the number of layers of $T$ and $N$ is its number of internal
gates.
\end{theorem}

\begin{theorem}[Switching failure probability le of half switching]
\lean{FABL.DepthCircuit.switchingFailureProbability_le_of_halfSwitching}
\statementsource{lmn-constant-depth-fourier}{
  pdf-63759bd909d5-p005-b004,
  pdf-63759bd909d5-p006-b001
}
\proofsource{Analysis-of-Boolean-Functions-by-Ryan-ODonnell}{
  pdf-ffbe9c7b3306-p104-b004,
  pdf-ffbe9c7b3306-p104-b005,
  pdf-ffbe9c7b3306-p105-b001,
  pdf-ffbe9c7b3306-p106-b002
}
\label{FABL.DepthCircuit.switchingFailureProbability_le_of_halfSwitching}
\leanok
\uses{FABL.CircuitGate.HasWidthLE, FABL.CircuitGate.dual, FABL.CircuitGate.hasWidthLE_mono, FABL.CircuitTail.internalNodeCount, FABL.CircuitTail.layerCount, FABL.CircuitTail.one_le_layerCount, FABL.CircuitTail.switchingFailureProbability_evalFunction_le_of_halfSwitching, FABL.DepthCircuit, FABL.DepthCircuit.depth, FABL.DepthCircuit.evalFunction_layer1, FABL.DepthCircuit.layer1Function, FABL.DepthCircuit.layer1Function_hasWidthLE, FABL.DepthCircuit.size, FABL.DepthCircuit.toBooleanFunction, FABL.DepthCircuit.width, FABL.HasHalfSwitchingBound, FABL.circuitCompressionRate, FABL.switchingFailureProbability}
\difficulty{4}
% dataset-pipeline: stub -- prose pending build_informal_statements.py
Statement pending regeneration.
\end{theorem}

\begin{definition}[Layer cutoff]
\lean{FABL.lmnLayerCutoff}
\label{FABL.lmnLayerCutoff}
\leanok
\uses{FABL.dnfWidthTruncationCutoff}
For a natural number $s$ and a real number $\varepsilon$, define the layer cutoff to be
the width-truncation cutoff at size parameter $s$ and tolerance $\varepsilon/2$.
\end{definition}

\begin{definition}[Output cutoff]
\lean{FABL.lmnOutputCutoff}
\label{FABL.lmnOutputCutoff}
\leanok
\uses{FABL.dnfWidthTruncationCutoff}
For a real number $\varepsilon$, define the output cutoff to be the width-truncation
cutoff at size parameter $1$ and tolerance $\varepsilon/2$.
\end{definition}

\begin{theorem}[Layer cutoff exponential bound]
\lean{FABL.mul_half_pow_lmnLayerCutoff_le}
\label{FABL.mul_half_pow_lmnLayerCutoff_le}
\leanok
\uses{FABL.lmnLayerCutoff, FABL.mul_inv_two_pow_dnfWidthTruncationCutoff_le}
\difficulty{3}
Let $s$ be a positive natural number and let $\varepsilon>0$. Define the layer cutoff
$L$ to be the least natural number greater than or equal to $\log_2(2s/\varepsilon)$.
Then $s(1/2)^L\le \varepsilon/2$.
\end{theorem}

\begin{theorem}[Output-Cutoff Power Bound]
\lean{FABL.half_pow_lmnOutputCutoff_le}
\label{FABL.half_pow_lmnOutputCutoff_le}
\leanok
\uses{FABL.lmnOutputCutoff, FABL.mul_inv_two_pow_dnfWidthTruncationCutoff_le}
\difficulty{3}
For every real number $\varepsilon>0$, let $d=\lceil\log_2(2/\varepsilon)\rceil$. Then
$(1/2)^d\le \varepsilon/2$.
\end{theorem}

\begin{theorem}[Positivity of the layer cutoff]
\lean{FABL.lmnLayerCutoff_pos}
\label{FABL.lmnLayerCutoff_pos}
\leanok
\uses{FABL.lmnLayerCutoff, FABL.mul_half_pow_lmnLayerCutoff_le}
\difficulty{4}
Let $s$ be a positive natural number and let $0<\varepsilon\leq1$. The least natural
number greater than or equal to $\log_2(2s/\varepsilon)$ is positive.
\end{theorem}

\begin{theorem}[Positivity of the output cutoff]
\lean{FABL.lmnOutputCutoff_pos}
\label{FABL.lmnOutputCutoff_pos}
\leanok
\uses{FABL.half_pow_lmnOutputCutoff_le, FABL.lmnOutputCutoff}
\difficulty{4}
Let $0<\varepsilon\leq1$. The least natural number greater than or equal to
$\log_2(2/\varepsilon)$ is positive.
\end{theorem}

\begin{theorem}[Switching bound for depth circuits]
\lean{FABL.lemma4_28_of_halfSwitching}
\proofsource{Analysis-of-Boolean-Functions-by-Ryan-ODonnell}{
  pdf-ffbe9c7b3306-p104-b003,
  pdf-ffbe9c7b3306-p104-b004,
  pdf-ffbe9c7b3306-p104-b005,
  pdf-ffbe9c7b3306-p105-b001,
  pdf-ffbe9c7b3306-p106-b002
}
\label{FABL.lemma4_28_of_halfSwitching}
\leanok
\uses{FABL.BooleanFunction, FABL.DepthCircuit.HasDepthCircuit, FABL.DepthCircuit.switchingFailureProbability_le_of_halfSwitching, FABL.DepthCircuit.toBooleanFunction, FABL.HasHalfSwitchingBound, FABL.circuitCompressionRate, FABL.half_pow_lmnOutputCutoff_le, FABL.lmnLayerCutoff, FABL.lmnLayerCutoff_pos, FABL.lmnOutputCutoff, FABL.mul_half_pow_lmnLayerCutoff_le, FABL.switchingFailureProbability}
\difficulty{4}
Assume that for every $q>0$, every sign-valued Boolean function on a finite cube having
a DNF or CNF of width at most $q$, and every integer $k\geq0$, the probability that a
random restriction with free-coordinate probability $1/(10q)$ leaves decision-tree depth
at least $k$ is at most $2^{-k}$. Let $f:\{-1,1\}^n\to\{-1,1\}$ be computed by a
depth-$d$ circuit of size at most $s$ and width at most $w$, where $s,w>0$, and let
$0<\varepsilon\leq1$. Set $L=\lceil\log_2(2s/\varepsilon)\rceil$ and
$K=\lceil\log_2(2/\varepsilon)\rceil$. Under a random restriction with free-coordinate
probability $\frac1{10w}\bigl(\frac1{10L}\bigr)^{\max(d-2,0)}$, the probability that the
restricted function has decision-tree depth at least $K$ is at most $\varepsilon$.
\end{theorem}

\begin{definition}[Width cutoff]
\lean{FABL.lmnWidthCutoff}
\label{FABL.lmnWidthCutoff}
\leanok
\uses{FABL.dnfWidthTruncationCutoff}
For a natural number $s$ and a real number $\varepsilon$, define the width cutoff to be
the width-truncation cutoff at size parameter $s$ and tolerance $\varepsilon/16$.
\end{definition}

\begin{definition}[Degree cutoff]
\lean{FABL.lmnDegreeCutoff}
\label{FABL.lmnDegreeCutoff}
\leanok
\uses{FABL.circuitCompressionRate, FABL.lmnLayerCutoff, FABL.lmnOutputCutoff, FABL.lmnWidthCutoff}
For natural numbers $d$ and $s$ and a real number $\varepsilon$, define the degree
cutoff to be the natural-number ceiling of three times the output cutoff at tolerance
$\varepsilon/12$, divided by the circuit compression rate at the width cutoff for $s$
and tolerance $\varepsilon$, and the layer cutoff for $s$ and tolerance
$\varepsilon/12$, raised to the power given by $d-2$ with truncated natural subtraction.
\end{definition}

\begin{theorem}[Boolean disagreement and relative Hamming distance]
\lean{FABL.uniformLpNorm_sub_toReal_sq_eq_four_mul_relativeHammingDist}
\label{FABL.uniformLpNorm_sub_toReal_sq_eq_four_mul_relativeHammingDist}
\leanok
\uses{FABL.BooleanFunction, FABL.BooleanFunction.toReal, FABL.relativeHammingDist, FABL.uniformInner, FABL.uniformLpNorm, FABL.uniformLpNorm_two_sq_eq_uniformInner, FABL.uniformProbability, FABL.uniformProbability_ne_eq_relativeHammingDist}
\difficulty{4}
Let $n\in\mathbb{N}$, and let $f,g:\{-1,1\}^n\to\{-1,1\}$ be Boolean-valued functions.
The squared uniform $2$-norm of the difference between their real-valued sign
representations equals four times their relative Hamming distance, namely four times the
number of inputs on which they disagree divided by $2^n$.
\end{theorem}

\begin{theorem}[Symmetry of relative Hamming distance]
\lean{FABL.relativeHammingDist_comm}
\label{FABL.relativeHammingDist_comm}
\leanok
\uses{FABL.BooleanFunction, FABL.relativeHammingDist}
\difficulty{3}
For any two Boolean functions on the sign cube of dimension $n$, their relative Hamming
distance is symmetric in the two functions.
\end{theorem}

\begin{theorem}[Positivity of the width cutoff]
\lean{FABL.lmnWidthCutoff_pos}
\label{FABL.lmnWidthCutoff_pos}
\leanok
\uses{FABL.lmnWidthCutoff, FABL.mul_inv_two_pow_dnfWidthTruncationCutoff_le}
\difficulty{4}
Let $s$ be a nonnegative integer and let $0<\varepsilon\le 1/2$, with $s>1$. The integer
$\lceil\log_2(16s/\varepsilon)\rceil$ is positive.
\end{theorem}

\begin{theorem}[Positivity of the restriction rate]
\lean{FABL.lmnRestrictionRate_pos}
\label{FABL.lmnRestrictionRate_pos}
\leanok
\uses{FABL.circuitCompressionRate, FABL.lmnLayerCutoff, FABL.lmnLayerCutoff_pos, FABL.lmnWidthCutoff, FABL.lmnWidthCutoff_pos, FABL.switchingLayerRate}
\difficulty{4}
Let $d,s$ be nonnegative integers and let $0<\varepsilon\le 1/2$, with $s>1$. Define
$W=\lceil\log_2(16s/\varepsilon)\rceil$ and $L=\lceil\log_2(24s/\varepsilon)\rceil$.
Then $\frac{1}{10W}\left(\frac{1}{10L}\right)^{d-2}>0$, where the exponent uses
truncated subtraction, so $d-2=\max(d-2,0)$.
\end{theorem}

\begin{theorem}[Truncation error at the width cutoff]
\lean{FABL.DepthCircuit.relativeHammingDist_truncateWidth_lmnWidthCutoff_le}
\proofsource{Analysis-of-Boolean-Functions-by-Ryan-ODonnell}{pdf-ffbe9c7b3306-p106-b005}
\label{FABL.DepthCircuit.relativeHammingDist_truncateWidth_lmnWidthCutoff_le}
\leanok
\uses{FABL.DepthCircuit, FABL.DepthCircuit.relativeHammingDist_truncateWidth_le, FABL.DepthCircuit.size, FABL.DepthCircuit.toBooleanFunction, FABL.DepthCircuit.truncateWidth, FABL.lmnWidthCutoff, FABL.mul_inv_two_pow_dnfWidthTruncationCutoff_le, FABL.relativeHammingDist}
\difficulty{4}
Let $n$ be a natural number, let $C$ be a depth circuit on the $n$-dimensional Boolean
cube, and let $s$ be a positive natural number with circuit size at most $s$. For every
$[?]$ positive real number $[?]$, truncating the widths of its first-layer terms at
$[?]$ gives a Boolean function whose relative Hamming distance from the function
computed by $C$ is at most $[?]$. Here the cutoff is $[?]$, the ceiling of the
base-$2$ logarithm of $16s/[?]$.
\end{theorem}

\begin{theorem}[Lmn theorem of half switching]
\lean{FABL.lmn_theorem_of_halfSwitching}
\statementsource{lmn-constant-depth-fourier}{
  pdf-63759bd909d5-p009-b001,
  pdf-63759bd909d5-p009-b002
}
\proofsource{Analysis-of-Boolean-Functions-by-Ryan-ODonnell}{
  pdf-ffbe9c7b3306-p106-b004,
  pdf-ffbe9c7b3306-p106-b005
}
\label{FABL.lmn_theorem_of_halfSwitching}
\leanok
\uses{FABL.BooleanFunction, FABL.BooleanFunction.toReal, FABL.DepthCircuit, FABL.DepthCircuit.HasDepthCircuit, FABL.DepthCircuit.relativeHammingDist_truncateWidth_lmnWidthCutoff_le, FABL.DepthCircuit.size_truncateWidth_le, FABL.DepthCircuit.toBooleanFunction, FABL.DepthCircuit.truncateWidth, FABL.DepthCircuit.width_truncateWidth_le, FABL.HasHalfSwitchingBound, FABL.IsFourierSpectrumConcentratedOn, FABL.IsFourierSpectrumConcentratedOn.transfer_of_uniformLpNorm_sub_sq_le, FABL.IsFourierSpectrumConcentratedUpTo, FABL.IsFourierSpectrumConcentratedUpTo.mono_cutoff, FABL.IsFourierSpectrumConcentratedUpTo.mono_error, FABL.circuitCompressionRate, FABL.circuitCompressionRate_le_one, FABL.isFourierSpectrumConcentratedOn_lowDegreeFourierFamily_iff, FABL.lemma4_21, FABL.lemma4_28_of_halfSwitching, FABL.lmnDegreeCutoff, FABL.lmnLayerCutoff, FABL.lmnLayerCutoff_pos, FABL.lmnOutputCutoff, FABL.lmnOutputCutoff_pos, FABL.lmnWidthCutoff, FABL.lmnWidthCutoff_pos, FABL.lowDegreeFourierFamily, FABL.relativeHammingDist, FABL.relativeHammingDist_comm, FABL.switchingFailureProbability, FABL.switchingLayerRate, FABL.uniformLpNorm, FABL.uniformLpNorm_sub_toReal_sq_eq_four_mul_relativeHammingDist}
\difficulty{5}
% dataset-pipeline: stub -- prose pending build_informal_statements.py
Statement pending regeneration.
\end{theorem}

\begin{definition}[Depth-size circuit class]
\lean{FABL.depthSizeCircuitClass}
\label{FABL.depthSizeCircuitClass}
\leanok
\uses{FABL.BooleanFunction, FABL.DepthCircuit.HasDepthCircuit}
Let $n$, $d$, and $s$ be natural numbers. The depth-size circuit class on $n$
coordinates consists of the Boolean functions for which there exists a width parameter
such that the function has a depth-circuit representation with parameters $d$, $s$, and
that width.
\end{definition}

\begin{definition}[Random example cost for a low-degree learner]
\lean{FABL.lowDegreeLearnerRandomExampleCost}
\statementsource{lmn-constant-depth-fourier}{pdf-63759bd909d5-p010-b002}
\label{FABL.lowDegreeLearnerRandomExampleCost}
\leanok
\uses{FABL.PositiveLearningParameter, FABL.finiteFamilySamplesPerCoefficient, FABL.lowDegreeFourierFamily, FABL.lowDegreeFourierFamily_nonempty}
For natural numbers $n$ and $k$ and a rational learning parameter $\varepsilon$ with
$0<\varepsilon\leq 1/2$, the random example cost is the cardinality of the family of
low-degree Fourier components on $n$ bits with cutoff $k$, multiplied by the number of
samples per coefficient for that family and parameter.
\end{definition}

\begin{definition}[Low-degree learner work cost]
\lean{FABL.lowDegreeLearnerWorkCost}
\statementsource{lmn-constant-depth-fourier}{
  pdf-63759bd909d5-p010-b001,
  pdf-63759bd909d5-p010-b002
}
\label{FABL.lowDegreeLearnerWorkCost}
\leanok
\uses{FABL.PositiveLearningParameter, FABL.lowDegreeLearnerRandomExampleCost}
For natural numbers $n$ and $k$ and a rational parameter $\tfrac12
obreak<
obreakrac12$?
\end{definition}

\begin{theorem}[Cost of the low-degree estimator]
\lean{FABL.lowDegreeFourierEstimatorProgram_cost_eq}
\label{FABL.lowDegreeFourierEstimatorProgram_cost_eq}
\leanok
\uses{FABL.BooleanFunction, FABL.LearningCost, FABL.LearningProgram.runWithCost, FABL.PositiveLearningParameter, FABL.SparseFourierHypothesis, FABL.finiteFamilyFourierEstimatorProgram_cost_eq, FABL.lowDegreeFourierEstimatorProgram, FABL.lowDegreeFourierFamily, FABL.lowDegreeFourierFamily_nonempty, FABL.lowDegreeLearnerRandomExampleCost, FABL.lowDegreeLearnerWorkCost}
\difficulty{1}
Let $n,k$ be natural numbers, let $\varepsilon$ be a positive learning parameter, and
let the target be a sign-valued Boolean function on the $n$-dimensional cube. Every
outcome in the support of the low-degree Fourier estimator's output law has cost
\[\bigl(R,0,R+R(n+1)\bigr),\qquad R=\text{the low-degree learner's random-example cost for }(n,k,\varepsilon).\]
\end{theorem}

\begin{theorem}[Cost of the low degree estimator]
\lean{FABL.lowDegreeFourierEstimatorProgram_cost_le_workBound}
\label{FABL.lowDegreeFourierEstimatorProgram_cost_le_workBound}
\leanok
\uses{FABL.BooleanFunction, FABL.LearningCost, FABL.LearningProgram.runWithCost, FABL.PositiveLearningParameter, FABL.SparseFourierHypothesis, FABL.lowDegreeFourierEstimatorProgram, FABL.lowDegreeFourierEstimatorProgram_cost_eq, FABL.lowDegreeLearnerWorkCost}
\difficulty{3}
Let $n,k$ be nonnegative integers, let $\varepsilon$ be a positive learning parameter,
and let $f$ be a Boolean function on the $n$-dimensional sign cube. If the estimator’s
work bound is at most a natural number $B$, then every outcome in the support of its run
on $f$ uses at most $B$ random examples, makes no membership queries, and has charged
work at most $B$.
\end{theorem}

\begin{theorem}[Learning circuits from low-degree concentration]
\lean{FABL.theorem4_31_of_lowDegreeConcentration}
\proofsource{lmn-constant-depth-fourier}{
  pdf-63759bd909d5-p010-b004,
  pdf-63759bd909d5-p010-b005,
  pdf-63759bd909d5-p010-b006,
  pdf-63759bd909d5-p010-b007,
  pdf-63759bd909d5-p010-b008,
  pdf-63759bd909d5-p010-b009,
  pdf-63759bd909d5-p011-b001
}
\proofsource{Analysis-of-Boolean-Functions-by-Ryan-ODonnell}{
  pdf-ffbe9c7b3306-p106-b008,
  pdf-ffbe9c7b3306-p106-b009
}
\label{FABL.theorem4_31_of_lowDegreeConcentration}
\leanok
\uses{FABL.BooleanFunction.toReal, FABL.IsFourierSpectrumConcentratedUpTo, FABL.LearningProgram.eventProbability, FABL.LearningProgram.runWithCost, FABL.PositiveLearningParameter, FABL.SparseFourierHypothesis.evaluate, FABL.depthSizeCircuitClass, FABL.lowDegreeFourierEstimatorProgram, FABL.lowDegreeFourierEstimatorProgram_cost_le_workBound, FABL.lowDegreeFourierEstimatorProgram_failureProbability_le_one_tenth, FABL.lowDegreeLearnerWorkCost, FABL.relativeHammingDist}
\difficulty{4}
Let $n,d,s,k,W$ be nonnegative integers, and let $\varepsilon\in\mathbb Q$ satisfy
$0<\varepsilon\le1/2$. Consider the Boolean functions on $\{-1,1\}^n$ computed by
circuits of depth $d$ and size at most $s$, with no bound imposed on their width.
Suppose every such function has total squared Fourier coefficient mass on sets of size
greater than $k$ at most $\varepsilon/2$, and suppose $W\ge R+R(n+1)$, where
$R=\binom{n}{\le k}\,M$, $\binom{n}{\le k}$ is the number of subsets of $[n]$ of size at
most $k$, and $M$ is the sample count prescribed by the Fourier estimator for accuracy
$\varepsilon/(2\binom{n}{\le k})$ and confidence parameter $1/(10\binom{n}{\le k})$.
Then the learner that estimates the Fourier coefficients on all those subsets and
returns the sign of the resulting polynomial has, for every function in the circuit
class, probability at most $1/10$ of producing a hypothesis whose relative Hamming
distance from the target exceeds $\varepsilon$. Every possible run uses at most $W$
labeled random examples, makes no membership queries, and incurs at most $W$ units of
work.
\end{theorem}

\begin{definition}[Circuit learning degree cutoff]
\lean{FABL.lmnCircuitLearningDegree}
\statementsource{lmn-constant-depth-fourier}{pdf-63759bd909d5-p010-b002}
\label{FABL.lmnCircuitLearningDegree}
\leanok
\uses{FABL.PositiveLearningParameter, FABL.lmnDegreeCutoff}
For natural numbers $d$ and $s$ and a rational number $\varepsilon$ satisfying
$0<\varepsilon\leq 1/2$, the circuit learning degree is the degree cutoff at parameters
$d$, $s$, and the real number $\varepsilon/2$.
\end{definition}

\begin{definition}[Work cost of the circuit learner]
\lean{FABL.lmnCircuitLearnerWorkCost}
\statementsource{lmn-constant-depth-fourier}{
  pdf-63759bd909d5-p010-b001,
  pdf-63759bd909d5-p010-b002
}
\label{FABL.lmnCircuitLearnerWorkCost}
\leanok
\uses{FABL.PositiveLearningParameter, FABL.lmnCircuitLearningDegree, FABL.lowDegreeLearnerWorkCost}
For natural numbers $n$, $d$, and $s$, and a rational learning parameter $\varepsilon$
satisfying $0<\varepsilon\leq 1/2$, define the circuit learner's work cost to be the
low-degree learner's work cost on $n$ inputs, with degree cutoff given by the circuit
learning degree determined by $d$, $s$, and $\varepsilon$.
\end{definition}

\begin{theorem}[Theorem4 31 of half switching]
\lean{FABL.theorem4_31_of_halfSwitching}
\statementsource{lmn-constant-depth-fourier}{
  pdf-63759bd909d5-p010-b004,
  pdf-63759bd909d5-p010-b009,
  pdf-63759bd909d5-p011-b001
}
\statementsource{Analysis-of-Boolean-Functions-by-Ryan-ODonnell}{
  pdf-ffbe9c7b3306-p106-b009,
  pdf-ffbe9c7b3306-p106-b004,
  pdf-ffbe9c7b3306-p106-b005
}
\label{FABL.theorem4_31_of_halfSwitching}
\leanok
\uses{FABL.HasHalfSwitchingBound, FABL.LearningProgram.eventProbability, FABL.LearningProgram.runWithCost, FABL.PositiveLearningParameter, FABL.SparseFourierHypothesis.evaluate, FABL.depthSizeCircuitClass, FABL.lmnCircuitLearnerWorkCost, FABL.lmnCircuitLearningDegree, FABL.lmn_theorem_of_halfSwitching, FABL.lowDegreeFourierEstimatorProgram, FABL.positiveLearningParameter_toReal_mem_Ioc, FABL.relativeHammingDist, FABL.theorem4_31_of_lowDegreeConcentration}
\difficulty{4}
% dataset-pipeline: stub -- prose pending build_informal_statements.py
Statement pending regeneration.
\end{theorem}

\begin{theorem}[Low-degree Fourier family size bound]
\lean{FABL.card_lmnCircuitLearningFamily_le}
\label{FABL.card_lmnCircuitLearningFamily_le}
\leanok
\uses{FABL.PositiveLearningParameter, FABL.card_lowDegreeFourierFamily_le, FABL.lmnCircuitLearningDegree, FABL.lowDegreeFourierFamily}
\difficulty{1}
Let $n,d,\mathrm{size}$ be natural numbers and let $\varepsilon\in\mathbb Q$ satisfy
$0<\varepsilon\leq 1/2$. The family of subsets of an $n$-element coordinate set having
size at most the prescribed circuit-learning degree for $d,\mathrm{size},\varepsilon$
has cardinality at most $(D+1)(n+1)^D$, where $D$ is that degree.
\end{theorem}

\begin{theorem}[Corollary4 32 of half switching]
\lean{FABL.DepthCircuit.corollary4_32_of_halfSwitching}
\statementsource{lmn-constant-depth-fourier}{pdf-63759bd909d5-p010-b004}
\statementsource{Analysis-of-Boolean-Functions-by-Ryan-ODonnell}{
  pdf-ffbe9c7b3306-p107-b002,
  pdf-ffbe9c7b3306-p107-b003
}
\label{FABL.DepthCircuit.corollary4_32_of_halfSwitching}
\leanok
\uses{FABL.DepthCircuit, FABL.DepthCircuit.HasDepthCircuit, FABL.DepthCircuit.depth, FABL.DepthCircuit.size, FABL.DepthCircuit.toBooleanFunction, FABL.DepthCircuit.width, FABL.HasHalfSwitchingBound, FABL.lmnDegreeCutoff, FABL.lmn_theorem_of_halfSwitching, FABL.parityAgreement_forces_concentration_cutoff, FABL.parityFunction, FABL.uniformProbability, FABL.uniformProbability_le_one}
\difficulty{5}
% dataset-pipeline: stub -- prose pending build_informal_statements.py
Statement pending regeneration.
\end{theorem}

\begin{theorem}[Hastad half switching bound for circuits]
\lean{FABL.hastadHalfSwitchingBoundForCircuits}
\statementsource{lmn-constant-depth-fourier}{
  pdf-63759bd909d5-p005-b004,
  pdf-63759bd909d5-p006-b001
}
\statementsource{Analysis-of-Boolean-Functions-by-Ryan-ODonnell}{
  pdf-ffbe9c7b3306-p104-b004,
  pdf-ffbe9c7b3306-p104-b005,
  pdf-ffbe9c7b3306-p105-b001,
  pdf-ffbe9c7b3306-p106-b002
}
\label{FABL.hastadHalfSwitchingBoundForCircuits}
\leanok
\uses{FABL.HasHalfSwitchingBound, FABL.hastadHalfSwitchingBound, FABL.switchingLayerRate}
\difficulty{1}
% dataset-pipeline: stub -- prose pending build_informal_statements.py
Statement pending regeneration.
\end{theorem}

\begin{theorem}[Lemma4 28]
\lean{FABL.lemma4_28}
\statementsource{Analysis-of-Boolean-Functions-by-Ryan-ODonnell}{
  pdf-ffbe9c7b3306-p104-b003,
  pdf-ffbe9c7b3306-p104-b004,
  pdf-ffbe9c7b3306-p104-b005,
  pdf-ffbe9c7b3306-p105-b001,
  pdf-ffbe9c7b3306-p106-b002
}
\label{FABL.lemma4_28}
\leanok
\uses{FABL.BooleanFunction, FABL.DepthCircuit.HasDepthCircuit, FABL.circuitCompressionRate, FABL.hastadHalfSwitchingBoundForCircuits, FABL.lemma4_28_of_halfSwitching, FABL.lmnLayerCutoff, FABL.lmnOutputCutoff, FABL.switchingFailureProbability}
\difficulty{1}
% dataset-pipeline: stub -- prose pending build_informal_statements.py
Statement pending regeneration.
\end{theorem}

\begin{theorem}[Lmn theorem]
\lean{FABL.lmn_theorem}
\statementsource{lmn-constant-depth-fourier}{
  pdf-63759bd909d5-p009-b001,
  pdf-63759bd909d5-p009-b002
}
\proofsource{Analysis-of-Boolean-Functions-by-Ryan-ODonnell}{
  pdf-ffbe9c7b3306-p106-b004,
  pdf-ffbe9c7b3306-p106-b005
}
\label{FABL.lmn_theorem}
\leanok
\uses{FABL.BooleanFunction, FABL.BooleanFunction.toReal, FABL.DepthCircuit.HasDepthCircuit, FABL.IsFourierSpectrumConcentratedUpTo, FABL.hastadHalfSwitchingBoundForCircuits, FABL.lmnDegreeCutoff, FABL.lmn_theorem_of_halfSwitching}
\difficulty{1}
% dataset-pipeline: stub -- prose pending build_informal_statements.py
Statement pending regeneration.
\end{theorem}

\begin{theorem}[Theorem4 31]
\lean{FABL.theorem4_31}
\statementsource{lmn-constant-depth-fourier}{
  pdf-63759bd909d5-p010-b004,
  pdf-63759bd909d5-p010-b009,
  pdf-63759bd909d5-p011-b001
}
\proofsource{Analysis-of-Boolean-Functions-by-Ryan-ODonnell}{pdf-ffbe9c7b3306-p106-b009}
\label{FABL.theorem4_31}
\leanok
\uses{FABL.LearningProgram.eventProbability, FABL.LearningProgram.runWithCost, FABL.PositiveLearningParameter, FABL.SparseFourierHypothesis.evaluate, FABL.depthSizeCircuitClass, FABL.hastadHalfSwitchingBoundForCircuits, FABL.lmnCircuitLearnerWorkCost, FABL.lmnCircuitLearningDegree, FABL.lowDegreeFourierEstimatorProgram, FABL.relativeHammingDist, FABL.theorem4_31_of_halfSwitching}
\difficulty{1}
% dataset-pipeline: stub -- prose pending build_informal_statements.py
Statement pending regeneration.
\end{theorem}

\begin{theorem}[Corollary4 32]
\lean{FABL.DepthCircuit.corollary4_32}
\statementsource{Analysis-of-Boolean-Functions-by-Ryan-ODonnell}{
  pdf-ffbe9c7b3306-p107-b002,
  pdf-ffbe9c7b3306-p107-b003
}
\label{FABL.DepthCircuit.corollary4_32}
\leanok
\uses{FABL.DepthCircuit, FABL.DepthCircuit.corollary4_32_of_halfSwitching, FABL.DepthCircuit.depth, FABL.DepthCircuit.size, FABL.DepthCircuit.toBooleanFunction, FABL.hastadHalfSwitchingBoundForCircuits, FABL.lmnDegreeCutoff, FABL.parityFunction, FABL.uniformProbability}
\difficulty{1}
% dataset-pipeline: stub -- prose pending build_informal_statements.py
Statement pending regeneration.
\end{theorem}
